% Source: https://artofproblemsolving.com/community/c3h1618182_complex_numbers
% Stable identity: aops-topic-1618182-q01
A number $z=x+yi$ has $r=|z|=\sqrt{x^2+y^2}$, and $\theta$ is the angle from the positive real axis to the point $(x,y)$. When the real part is $0$ the formula $\tan\theta=\dfrac {y}{x}$ cannot be used, but the point $(0,y)$ lies on the imaginary axis, so the angle can be read off directly.
\begin{itemize}
\item $4i$ is the point $(0,4)$ on the positive imaginary axis: $r=4$, $\theta=90^\circ$, so $4i=4\left(\cos90^\circ+i\sin90^\circ\right)$.
\item $-3i$ is the point $(0,-3)$ on the negative imaginary axis: $r=3$, $\theta=270^\circ$, so $-3i=3\left(\cos270^\circ+i\sin270^\circ\right)$.
\end{itemize}
In general
\[
\boxed{bi=b\left(\cos90^\circ+i\sin90^\circ\right)\text{ if }b>0,\qquad bi=|b|\left(\cos270^\circ+i\sin270^\circ\right)\text{ if }b<0}.
\]
(Check: $\cos90^\circ=0$, $\sin90^\circ=1$; $\cos270^\circ=0$, $\sin270^\circ=-1$.)
